\documentclass[troisieme]{fiche}
\metadonnees{4}{Jouer aux pauvres}{Bloc A — Raconter}

\begin{document}
\entete

\objectifs{%
\textbf{Langue} : le present perfect ; \emph{for} et \emph{since}.\par
\textbf{Texte} : E. Nesbit, \emph{The Railway Children} (1906), ch. \textsc{i}.\par
\textbf{Durée} : 1 h — exercices 20 min, texte et questions 25 min, version 15 min.}

%% ===================================================================
\exercice[12 min]{Present perfect ou prétérit}

\rappel{\textbf{Rappel.} Le present perfect (\emph{have} + participe passé) relie le passé à
maintenant : le résultat compte encore. Le prétérit coupe le lien : il faut un moment précis.}

\consigne{Mettez le verbe entre parenthèses au temps qui convient.}

\begin{anglais}
\begin{enumerate}[label=\arabic*.,itemsep=2.5mm,leftmargin=*]
  \item We \emph{(live)} \trou{have lived} in this house since 1900.
  \item They \emph{(leave)} \trou{left} London last week.
  \item Mother \emph{(be)} \trou{has been} ill for two days.
  \item The doctor \emph{(come)} \trou{came} on Tuesday morning.
  \item I \emph{(never see)} \trou{have never seen} such a big house.
  \item Peter \emph{(get)} \trou{got} a model engine for his tenth birthday.
  \item \emph{(you / ever / travel)} \trou{Have you ever travelled} by train at night?
  \item The men \emph{(just take)} \trou{have just taken} the furniture away.
  \item When \emph{(they / pack)} \trou{did they pack} the crockery?
  \item Roberta \emph{(not forget)} \trou{has not forgotten} her mother's face.
\end{enumerate}
\end{anglais}

\note{Present perfect : 1, 3, 5, 7, 8, 10. Prétérit : 2, 4, 6, 9.}

\note[Le repère décide]{\emph{Last week}, \emph{on Tuesday}, \emph{when\ldots?} donnent un
moment précis : prétérit. \emph{Since}, \emph{for}, \emph{ever}, \emph{never}, \emph{just}
relient à maintenant : present perfect. \emph{When have you arrived?} est impossible en
anglais.}

%% ===================================================================
\exercice[8 min]{\emph{For} ou \emph{since}}

\rappel{\textbf{Rappel.} \emph{For} annonce une durée, \emph{since} un point de départ. Les
deux s'emploient avec le present perfect.}

\consigne{Complétez avec \emph{for} ou \emph{since}.}

\begin{anglais}
\begin{enumerate}[label=\arabic*.,itemsep=2.5mm,leftmargin=*]
  \item They have lived in the Red Villa \trou{for} ten years.
  \item Mother has been in bed \trou{since} Monday.
  \item We have not seen Father \trou{for} three days.
  \item The house has been empty \trou{since} the beginning of the month.
  \item I have known her \trou{since} we were children.
  \item It has been raining \trou{for} hours.
  \item Nothing has changed \trou{since} last summer.
  \item She has been a teacher \trou{for} twenty years.
\end{enumerate}
\end{anglais}

\note[Deux remarques]{\emph{Since} peut être suivi d'un groupe nominal (\emph{since Monday})
ou d'une proposition au prétérit (\emph{since we were children}). Attention au français :
\og depuis dix ans \fg{} se dit \emph{for ten years} — \emph{since} ne traduit
\og depuis \fg{} que devant une date ou un événement.}

%% ===================================================================
\newpage
\section{Texte}

\noindent\small\textit{Trois enfants vivent une vie confortable dans la banlieue de Londres.
Un soir, deux hommes sont venus chercher leur père, et personne ne leur a dit
pourquoi.}\normalsize

\begin{texte}
Then came the time when Mother came home and went to bed and stayed there two days and the
Doctor came, and the children crept \gl[wretchedly]{wretchedly}{misérablement} about the house
and wondered if the world was coming to an end.

Mother came down one morning to breakfast, very pale and with lines on her face that used not
to be there. And she smiled, as well as she could, and said:

``Now, my \gl[a pet]{pets}{mes chéris}, everything is settled. We're going to leave this
house, and go and live in the country. Such a \gl[ducky]{ducky}{adorable} dear little white
house. I know you'll love it.''

A whirling week of packing followed --- not just packing clothes, like when you go to the
seaside, but packing chairs and tables, covering their tops with
\gl[sacking]{sacking}{de la toile de jute} and their legs with straw.

All sorts of things were packed that you don't pack when you go to the seaside.
\gl[crockery]{Crockery}{la vaisselle}, blankets, \gl[a candlestick]{candlesticks}{un
bougeoir}, carpets, \gl[a bedstead]{bedsteads}{un châlit, un cadre de lit},
\gl[a saucepan]{saucepans}{une casserole}, and even \gl[a fender]{fenders}{un pare-feu} and
fire-irons.

The house was like a furniture \gl[a warehouse]{warehouse}{un entrepôt}. I think the children
enjoyed it very much. Mother was very busy, but not too busy now to talk to them, and read to
them, and even to make a bit of poetry for Phyllis to cheer her up when she fell down with a
\gl[a screwdriver]{screwdriver}{un tournevis} and ran it into her hand.

``Aren't you going to pack this, Mother?'' Roberta asked, pointing to the beautiful cabinet
\gl[to inlay, inlaid]{inlaid}{incruster} with red \gl[turtleshell]{turtleshell}{l'écaille de
tortue} and brass.

``We can't take everything,'' said Mother.

``But we seem to be taking all the ugly things,'' said Roberta.

``We're taking the useful ones,'' said Mother; ``we've got to play at being Poor for a bit,
my \gl[a chickabiddy]{chickabiddy}{mon poussin (mot tendre)}.''

[\ldots]

``I say, this is \gl[larks]{larks}{c'est une rigolade},'' said Peter, wriggling joyously, as
Mother \gl[to tuck in]{tucked}{border (dans un lit)} him up. ``I do like moving! I wish we
moved once a month.''

Mother laughed.

``I don't!'' she said. ``Good night, Peterkin.''

As she turned away Roberta saw her face. She never forgot it.
\end{texte}
\source{E. Nesbit, \emph{The Railway Children}, Londres, Wells Gardner, 1906,
ch.~\textsc{i}.}

\partiel[15 min]{Questions}
\consigne{Répondez aux deux premières questions en français, aux deux dernières en anglais.}

\begin{enumerate}[label=\arabic*.,itemsep=2.5mm,leftmargin=*]
  \item Qu'est-ce que les enfants comprennent de la situation, au début du texte ?
    \lignes{2}
    \corr{Rien de précis. Leur mère est restée deux jours au lit, le médecin est venu, et ils
    rôdent dans la maison en se demandant si le monde va finir. Personne ne leur explique
    quoi que ce soit.}
  \item Qu'est-ce que Roberta remarque pendant les préparatifs, et que lui répond sa mère ?
    \lignes{2}
    \corr{Qu'on n'emporte que les objets laids, et qu'on laisse le beau meuble incrusté
    d'écaille. Sa mère répond qu'on emporte les objets utiles, et qu'ils vont devoir
    \og jouer aux pauvres \fg{} quelque temps.}
  \item \en{Why does Peter say he likes moving, and why does his mother answer ``I don't''?}
    \lignes{3}
    \corr{For Peter it is an adventure: the house upside down, sleeping on the sofa. His
    mother knows they are leaving because they are ruined, and that the father is not coming
    back. The same event is a game for one and a disaster for the other.}
  \item \en{``We've got to play at being Poor for a bit.'' Find another sentence in the text
    where a verb is contracted with \emph{have} or \emph{be}.}
    \lignes{2}
    \corr{\emph{We're going to leave this house}, \emph{you'll love it}, \emph{Aren't you
    going to pack this?}, \emph{We can't take everything}, \emph{I don't!} La langue parlée
    contracte ; le récit, lui, ne contracte pas.}
\end{enumerate}

\note[Le detail qui compte]{\emph{Lines on her face that used not to be there} : des rides
qui n'y étaient pas avant. Et la dernière phrase, \emph{As she turned away Roberta saw her
face. She never forgot it.} : la mère ne sourit que devant les enfants. Tout le chapitre tient
dans ces deux notations.}

%% ===================================================================
\section{Version}
\consigne{Traduisez le passage en français. Ce sont les premières lignes du roman.}

\begin{version}
They were not railway children to begin with. I don't suppose they had ever thought about
railways except as a means of getting to the Pantomime, the Zoological Gardens, and Madame
Tussaud's. They were just ordinary \gl[suburban]{suburban}{de banlieue} children, and they
lived with their Father and Mother in an ordinary red-brick-fronted
\gl[a villa]{villa}{un pavillon}, with coloured glass in the front door, a
\gl[tiled]{tiled}{carrelé} passage that was called a hall, a bath-room with hot and cold
water, electric bells, French windows, and a good deal of white paint, and `every modern
\gl[a convenience]{convenience}{un confort, un équipement}', as the house-agents say.
\end{version}
\source{\emph{Ibid.}}

\lignes{12}

\corr{\textbf{Proposition de traduction.} Ce n'étaient pas, au départ, des enfants du chemin
de fer. Je suppose qu'ils n'avaient jamais pensé aux trains autrement que comme un moyen
d'aller à la pantomime, au jardin zoologique et chez Madame Tussaud. C'étaient de simples
enfants de banlieue, et ils vivaient avec leur père et leur mère dans un pavillon ordinaire à
façade de brique rouge, avec des vitraux de couleur dans la porte d'entrée, un couloir carrelé
qu'on appelait le vestibule, une salle de bains avec eau chaude et eau froide, des sonnettes
électriques, des portes-fenêtres, beaucoup de peinture blanche, et \og tout le confort
moderne \fg, comme disent les agents immobiliers.}

\note[Choix de traduction 1]{\emph{to begin with} : \og pour commencer \fg, \og au départ \fg.
Locution figée, à ne pas découper.}

\note[Choix de traduction 2]{\emph{a means of getting to} : \emph{means} est invariable et
toujours écrit avec un \emph{s}. Après une préposition, l'anglais met le gérondif :
\emph{of getting} $\rightarrow$ \og d'aller \fg.}

\note[Choix de traduction 3]{\emph{a tiled passage that was called a hall} : le narrateur se
moque gentiment ; le mot \emph{hall} est trop grand pour la chose. Garder l'ironie en
traduisant \og qu'on appelait le vestibule \fg{} plutôt que \og appelé couloir \fg.}

%% ===================================================================
\partiel{Lexique à mémoriser}
\begin{multicols}{2}\small
\begin{itemize}[label=--,itemsep=0.8mm,leftmargin=*]
  \item \textbf{to creep} — se glisser, ramper
  \item \textbf{pale} — pâle
  \item \textbf{settled} — réglé, décidé
  \item \textbf{to pack} — faire les bagages, emballer
  \item \textbf{crockery} — la vaisselle
  \item \textbf{a blanket} — une couverture
  \item \textbf{a carpet} — un tapis
  \item \textbf{furniture} — les meubles \textit{(indénombrable)}
  \item \textbf{busy} — occupé
  \item \textbf{to cheer up} — remonter le moral
  \item \textbf{ugly} — laid
  \item \textbf{useful} — utile
\end{itemize}
\end{multicols}

\end{document}
